\section{Dónde encaja este trabajo}
\label{sec:posicionamiento}

Antes de la formulación conviene situar el problema, porque la automatización de
almacenes tiene un vocabulario de decisiones consolidado y no todas las
decisiones que enumera son objeto de este trabajo. Declarar cuáles sí y cuáles
no evita que el lector suponga cobertura donde no la hay.

\subsection{Los tres niveles de decisión}

La literatura de automatización de almacenes ordena las decisiones en tres
horizontes \cite{yildirim2023warehouse}: \emph{estratégico} —selección de
sistema, distribución en planta, trazado de pasillos, ubicación de puntos de
recogida y entrega, ubicación de estaciones de carga—; \emph{táctico}
—asignación de almacenamiento, gestión de pedidos, dimensionado de flota,
mantenimiento y gestión de fallos, gestión energética—; y \emph{operativo}
—asignación de tareas, planificación de rutas, resolución de bloqueos y
prevención de conflictos.

\begin{table}[!ht]
\centering\small
\caption{Decisiones industriales y resultados de este trabajo.}
\label{tab:posicionamiento}
\begin{tabularx}{\textwidth}{@{}p{2.0cm}p{4.0cm}X@{}}
\toprule
Nivel & Decisión & Resultado que la aborda \\
\midrule
Estratégico & trazado de pasillos y anchos
  & \cref{pr:compatibilidad}: la cota de grado del árbol de conectividad depende
    de $d_{\min}/R^{\mathrm{com}}$; \cref{obs:precio-astop}: el ancho del pasillo
    fija la velocidad admisible \\
Estratégico & emplazamiento de la reserva
  & \cref{th:potencia} y \cref{pr:ascenso}: partición por diagrama de potencia
    y reubicación distribuida que maximiza la cobertura de rescate \\
Estratégico & ubicación de estaciones de carga
  & \emph{no abordada} \\
Estratégico & selección de sistema, ergonomía, interacción humana
  & \emph{no abordadas} \\
\midrule
Táctico & gestión de fallos
  & \cref{th:subprovision} y \cref{th:prima}: nivel de reserva y su mecanismo \\
Táctico & gestión energética
  & \cref{pr:caudal}: caudal sostenido con energía como restricción lineal;
    \cref{obs:media-energia}: por qué la media no basta;
    \cref{obs:recarga}: programación de recarga con capacidad de red \\
Táctico & gestión de pedidos
  & \cref{co:admision}: control de admisión que preserva la estabilidad de
    precios \\
Táctico & dimensionado de flota
  & \emph{parcial}: \cref{pr:caudal} acota el caudal de una flota dada, no
    resuelve su tamaño óptimo \\
\midrule
Operativo & asignación de tareas
  & núcleo del trabajo: \cref{th:potencial}, \cref{co:terminacion},
    \cref{th:noequiv}, \cref{pr:jerarquia} \\
Operativo & planificación de rutas
  & \cref{th:multiclase}: peaje marginal por clase de envolvente \\
Operativo & bloqueos y conflictos
  & reserva vecinal con garantía de espera acotada;
    \cref{obs:posicionamiento-almacen} \\
\bottomrule
\end{tabularx}
\end{table}

\begin{observacion}[Concentración y sus huecos]
\label{obs:concentracion}
El trabajo es denso en el nivel operativo, aporta al táctico y alcanza el
estratégico por dos vías: la geométrica —anchos de pasillo— y la de cobertura
—emplazamiento de la reserva, \cref{sec:cobertura}—. Tres decisiones
industriales quedan explícitamente fuera: la selección de sistema, la
interacción con operarios y la ubicación de estaciones de carga. Ninguna es un olvido: todas dependen de
factores —coste, normativa, ergonomía— ajenos al formalismo empleado aquí.
\end{observacion}

\subsection{Qué distingue este problema del transporte por robot único}

La literatura de robótica cooperativa clasifica los sistemas multi-agente por
tres ejes \cite{gorrostieta2019applications}: tipo de agentes (homogéneos u
heterogéneos), arquitectura de control (reactiva, deliberativa o híbrida) y
comunicación (implícita o explícita). Este trabajo se sitúa en heterogéneos,
híbrida y explícita, con la precisión de que la comunicación explícita está
acotada a $n$ saltos y su implementabilidad se verifica con
\cref{pr:implementable}.

Lo que separa este problema de la asignación de tareas convencional no es la
heterogeneidad —que es común— sino que la ejecución exige \textbf{contacto físico simultáneo de varios agentes sobre un mismo cuerpo}. De ahí las tres
consecuencias que estructuran el trabajo:

\begin{enumerate}[leftmargin=1.9em,topsep=0.2em,itemsep=0.3em]
  \item la coalición debe certificar que puede producir el \emph{wrench}
  requerido, y ese certificado tiene fronteras (\cref{obs:fronteras});
  \item una vez acoplada, deja de ser un conjunto de puntos y pasa a ser un
  cuerpo con envolvente, lo que cambia su consumo de espacio
  (\cref{def:clase});
  \item su capacidad se suma pero su movilidad se intersecta
  (\cref{co:agarre}), de modo que coaliciones mayores no son necesariamente
  mejores.
\end{enumerate}

\begin{observacion}[La frontera con el control de formación]
\label{obs:frontera-formacion}
El control de formación cooperativo resuelve que un conjunto de robots mantenga
una geometría relativa mientras sigue una referencia. Ese problema no incluye la
pregunta que aquí es central: \emph{qué esfuerzo aporta cada miembro y si la
suma realiza la demanda}. La revisión de la bibliografía de robótica móvil y
control cooperativo consultada confirma esa separación
(\cref{obs:biblioteca-fisica}): ninguna de esas obras formula el reparto por
\emph{wrench} resultante con unilateralidad y conos de fricción. La formación
es condición necesaria del transporte cooperativo; no es su formulación.
\end{observacion}

\subsection{La arquitectura, nombrada}

El diseño de este trabajo sigue una metodología establecida sin haberla
declarado. Conviene nombrarla, porque el vocabulario aclara qué hace cada capa
\cite{ren2008distributed}: (1) fijar el objetivo de cooperación y sus
restricciones; (2) identificar la \emph{variable de coordinación} —la
información mínima que cada vehículo necesita conocer para coordinarse— y la
\emph{función de coordinación} que la liga al objetivo; (3) resolver el problema
suponiendo conocimiento global; (4) construir consenso para que todos dispongan
de esa información pese a una red imperfecta.

\begin{table}[!ht]
\centering\small
\caption{La metodología de \cite{ren2008distributed} y su realización aquí.}
\label{tab:metodologia}
\begin{tabularx}{\textwidth}{@{}p{3.4cm}X@{}}
\toprule
Paso & Realización en este trabajo \\
\midrule
Objetivo y restricciones
  & realizar $\Wdem_\ell$ con esfuerzos admisibles; restricciones de
    unilateralidad, cono de fricción, barreras y cotas de actuador \\
Variable de coordinación
  & el vector de precios $p_\ell$ del ajuste \eqref{eq:tatonnement} \\
Función de coordinación
  & el potencial exacto de \cref{th:potencial}, que liga la variable con el
    objetivo \\
Esquema centralizado
  & el programa convexo cuyo dual son esos precios \\
Construcción de consenso
  & \eqref{eq:tatonnement} difundido sobre el grafo de
    \cref{th:jerarquia-grafos} \\
\bottomrule
\end{tabularx}
\end{table}

\begin{observacion}[Por qué el precio es la variable de coordinación correcta]
\label{obs:variable-coordinacion}
La elección no es cosmética. Si la coalición acordase directamente el
\emph{reparto} —quién aporta cuánto— tendría que consensuar $\lvert\Ccal_\ell
\rvert$ magnitudes, y la carga de comunicación crecería con el tamaño de la
coalición. Acordando el \emph{precio} se consensúan tres magnitudes, la dimensión
del \emph{wrench} en el plano, sea cual sea el número de miembros: $2$, $8$ o
$32$ AMR acuerdan siempre tres números, y cada uno deriva su esfuerzo
localmente.

Esa es la razón estructural de que la arquitectura escale, y explica por qué
\cref{obs:campo-medio} puede afirmar que crecer la flota no la degrada. Es
también la forma precisa de decir en qué sentido este trabajo es dual: se
consensúa el multiplicador, no la asignación.
\end{observacion}

\begin{observacion}[Los cuatro acoplamientos, todos presentes]
\label{obs:acoplamientos}
La misma fuente clasifica el acoplamiento entre vehículos en cuatro tipos: de
objetivo, local, pleno y dinámico. La mayoría de los problemas de control
cooperativo exhiben uno o dos. Este los tiene los cuatro: de objetivo, porque la
demanda de \emph{wrench} es común; local, por las barreras entre vecinos; pleno,
por la congestión de \cref{th:multiclase}, que acopla a todos a través de los
arcos; y dinámico, porque en \textsf{TRANS} los AMR comparten un cuerpo rígido y
sus dinámicas quedan ligadas por \cref{eq:carga-dinamica}. Enumerarlo así es la
forma más económica de decir por qué el problema no se reduce a ninguna de sus
partes.
\end{observacion}

\begin{observacion}[Dos hechos estructurales del consenso de segundo orden]
\label{obs:segundo-orden}
Dado que los AMR son sistemas de segundo orden, dos resultados de
\cite{ren2008distributed} acotan lo que la capa de consenso puede y no puede
hacer, y ambos se han comprobado.

El primero: el consenso de doble integrador \emph{requiere} acoplamiento en la
velocidad relativa. Con él, la dispersión cae a $3\times10^{-14}$; acoplando
solo posiciones, el sistema oscila indefinidamente y la dispersión se estanca en
$1{,}95$. No es una cuestión de ganancias: el protocolo sin término de velocidad
no converge. Esto sostiene la exigencia de \cref{obs:hibrido-restante} de un
observador cuando la velocidad del vecino no se mide directamente.

El segundo, tranquilizador: el consenso sobrevive a la saturación de las
entradas. Con límites de $1$, $0{,}2$ y hasta $0{,}05$ el sistema sigue
convergiendo —dispersiones de $10^{-14}$ a $10^{-7}$—, aunque más despacio. Las
cotas de actuador de \cref{pr:cota} no comprometen, por tanto, la capa de
acuerdo; solo la retardan.
\end{observacion}

\subsection{Frente a la literatura de asignación en almacenes}

La revisión sistemática de \cite{yildirim2023warehouse} tabula los enfoques
publicados de asignación de tareas y de planificación de rutas en almacenes
robotizados. Conviene agruparlos para ver qué familia falta.

\begin{table}[!ht]
\centering\small
\caption{Familias de asignación de tareas en la literatura de almacenes, según
la tabulación de \cite{yildirim2023warehouse}, y su relación con este trabajo.}
\label{tab:asignacion-lit}
\begin{tabularx}{\textwidth}{@{}p{2.7cm}p{4.5cm}X@{}}
\toprule
Familia & Enfoques tabulados & Relación \\
\midrule
Reglas de despacho
  & vehículo más cercano, menor tiempo de respuesta, plazo, prioridad,
    uno de $N$
  & sin garantía de optimalidad; no comprueban factibilidad física \\
Metaheurísticas
  & algoritmo genético, recocido simulado, búsqueda por entornos amplia
  & óptimo no certificado; coste de cómputo no acotado por intervalo \\
Aprendizaje
  & \emph{deep Q-network}, redes neuronales, aprendizaje por refuerzo
  & sin certificado por instancia \\
Optimización exacta
  & programación lineal, entera y mixta, flujo de coste mínimo con conflictos
  & óptimo, pero centralizado y sin modelo de contacto \\
Mercado
  & subasta, negociación por red de contratos
  & la familia más próxima; asigna tarea a \emph{un} robot, no reparte
    esfuerzo entre varios \\
\midrule
\emph{Este trabajo}
  & juego de potencial exacto con certificado de \emph{wrench}
  & \cref{th:potencial}, \cref{th:span}, \cref{th:robusto} \\
\bottomrule
\end{tabularx}
\end{table}

\begin{observacion}[Qué no aparece en esa tabulación]
\label{obs:hueco-literatura}
Ninguno de los enfoques tabulados hace las dos cosas que aquí son centrales.
Primera: \emph{condicionar la asignación a un certificado de capacidad física}.
Las reglas y metaheurísticas asignan por distancia, plazo o coste, sin comprobar
que el conjunto asignado pueda producir el esfuerzo requerido; en un problema de
transporte cooperativo eso permite formar coaliciones que no pueden mover la
carga. Segunda: \emph{disponer de una función potencial}, que es lo que convierte
la revisión asíncrona de decisiones en un proceso con terminación garantizada
(\cref{co:terminacion}) en lugar de en una heurística evaluada empíricamente.

La familia de mercado —subastas y red de contratos— es la más cercana en
espíritu, pero su unidad de asignación es la pareja robot--tarea. El objeto que
aquí se subasta no es una tarea sino una \emph{componente de esfuerzo} sobre un
cuerpo compartido, y su precio se ajusta por un proceso de \emph{tâtonnement}
cuyo punto de reposo es el reparto que realiza la demanda.
\end{observacion}

\begin{observacion}[Planificación de rutas: todos planifican para puntos]
\label{obs:rutas-lit}
La tabulación de rutas de \cite{yildirim2023warehouse} recoge esencialmente dos
familias: variantes de $A^{\!*}$ —con retícula de estados, cooperativo,
jerárquico con ventana, con reserva de ruta y tiempos de giro— y variantes de
Dijkstra. Ambas planifican para vehículos tratados como puntos o como discos de
tamaño fijo. Una coalición acoplada no es ninguna de las dos cosas: su
envolvente depende de qué AMR la componen y de cómo se disponen, y cambia al
recomponerse. Esa es la razón de introducir clases de envolvente
(\cref{def:clase}) y el peaje por clase de \cref{th:multiclase}, que ninguna de
las dos familias contempla.
\end{observacion}

\subsection{Los cuatro conflictos y qué hace el mecanismo con cada uno}

La misma obra clasifica los conflictos entre robots en cuatro tipos, con tres
respuestas típicas —esperar, replanificar o apartarse— y zonificación estática o
dinámica como prevención.

\begin{table}[!ht]
\centering\small
\caption{Los cuatro tipos de conflicto de \cite{yildirim2023warehouse} y el
resultado de este trabajo que los cubre.}
\label{tab:conflictos}
\begin{tabularx}{\textwidth}{@{}p{2.5cm}p{4.3cm}X@{}}
\toprule
Tipo & Descripción & Tratamiento aquí \\
\midrule
Cruce & dos robots alcanzan la intersección a la vez
  & testigo de reserva vecinal con desempate lexicográfico; espera acotada \\
Frontal & dos robots se encuentran en sentidos opuestos
  & reserva más peaje por clase: la envolvente ancha paga el espacio que
    consume (\cref{th:multiclase}) \\
Persecución & un robot alcanza a otro más lento en la misma ruta
  & es el caso que el peaje internaliza de forma nativa: una coalición es una
    clase lenta y ancha, y su peaje $p_c\,y_e\,t_e'(y_e)$ cobra exactamente la
    demora que impone a los demás \\
Ocupación & un robot pretende ir donde otro está detenido
  & la detención en \textsf{DELIV} termina por \cref{pr:extincion}; la
    detención por parálisis de autoridad se diagnostica con
    \cref{th:degradacion} \\
\bottomrule
\end{tabularx}
\end{table}

\begin{observacion}[El conflicto de persecución es el argumento del peaje]
\label{obs:persecucion}
De los cuatro, el de persecución es el que mejor justifica la construcción de
\cref{sec:multiclase}. Las respuestas tabuladas en la literatura —esperar,
replanificar, apartarse— tratan el encuentro como un evento a resolver. El peaje
marginal lo trata como un coste a internalizar antes de que ocurra: si una
coalición ancha y lenta va a ocupar un arco, paga por la demora que impondrá, y
esa señal cambia la ruta que elige. Es la diferencia entre gestionar el conflicto
y no generarlo, y es exactamente el desplazamiento hacia enfoques proactivos que
la propia obra señala como frontera de la práctica.
\end{observacion}

\begin{observacion}[Indicadores: qué mide este trabajo y qué no]
\label{obs:kpis}
La misma fuente separa los indicadores de almacén —financieros, de cliente, de
proceso— de los específicos del sistema robótico: caudal de sistema, utilización
de robots, movimientos de estantería, eficiencia de recursos. De estos, el marco
de \cref{pr:caudal} produce directamente el \emph{caudal sostenido} y la
\emph{utilización}, ambos como salidas del programa lineal de medidas de
ocupación, y admite la energía como restricción explícita. No produce los
indicadores financieros ni los de cliente, que dependen de costes y de
compromisos de servicio ajenos al modelo. Declarar esa frontera importa porque
el caudal es el indicador que la literatura consultada considera el principal, y
es precisamente el que aquí se acota.
\end{observacion}

\subsection{Cinco referentes recientes y dónde queda este trabajo}

Hay cinco trabajos de 2024--2026 que resuelven partes de lo que aquí se
construye, y conviene decir frente a cada uno qué se comparte y qué no
\cite{mayorga2026compendio}.

\begin{center}\small
\begin{tabular}{@{}p{3.6cm}p{5.2cm}p{5.6cm}@{}}
\toprule
Referente & Qué resuelve & Qué queda fuera de su modelo y dentro de este \\
\midrule
DisCo \cite{shorinwa2024disco}: contacto implícito y coordinación distribuida
  & fuerzas y cambios de contacto optimizados en problemas locales con
    variables duales y comunicación vecinal; aproximarse, tocar y transportar
    sin máquina de estados
  & flota heterogénea en operación continua, información de $n$ saltos,
    apoyo inferior en \textsf{cargo}, certificados de ejecución bajo una
    física declarada (\cref{th:cerco}, \cref{pr:disco-friccion}); es la
    referencia central, no secundaria \\
MPC descentralizado para dos AMR diferenciales \cite{muhammed2026decentralized}
  & dos AMR acoplados físicamente, problemas locales con restricciones de
    coordinación y obstáculos, validación en hardware
  & referencias tomadas de un camino geométrico previo; aquí participación y
    movimiento responden a la misma función de valor
    (\cref{def:continuaciones}); su validación física es la que este
    documento no tiene \\
NMPC distribuido por ADMM con \emph{wrench} y barreras de orden superior \cite{sambhus2026distributed}
  & dinámica de carga, acoplamientos holónomos, consenso de trayectorias y
    \emph{wrench}, barreras para agentes y carga; simulación y experimentos
  & reclutamiento variable, apoyo inferior, decisiones mecánicas marginales
    (\cref{th:actuador-adicional}) y ausencia de modos nominales; no se
    reclama prioridad en coordinar \emph{wrench}, contacto y seguridad de
    forma distribuida \\
CT-TAPF / CT-TCBS \cite{zhou2026cttcbs}
  & formación de equipos, asignación y caminos sin colisión integrados, con
    solucionador óptimo y variantes subóptimas
  & su optimalidad es respecto de un problema discreto; no se traslada a
    tiempo o energía de la planta de contacto y par
    (\cref{def:jerarquia-optimos}) \\
DeReCo \cite{shibata2026dereco}
  & aprendizaje multiagente con representación y coordinación separadas,
    ejecución descentralizada desde observaciones locales, objetos no vistos
  & adaptación por modelo y precios frente a adaptación aprendida: aquí no se
    entrena una política por campaña pero se depende del modelo
    (\cref{obs:parametros-compartidos}); no se supone mejor generalización \\
\bottomrule
\end{tabular}
\end{center}

\begin{observacion}[El régimen donde este enfoque puede ganar, y dónde no]
\label{obs:regimen-ventaja}
Con una carga, equipo fijo y buena comunicación, un programa cuadrático
central con impedancia bien ajustada es difícil de superar y el mercado añade
negociación innecesaria; con obstáculos que exigen anticipar varios giros, el
horizonte de un NMPC evita trampas del campo reactivo; con asignación
combinatoria ajustada, una relajación continua no resuelve la atomicidad
(\cref{th:mohr}); con propiedades físicas parcialmente desconocidas, los
métodos aprendidos pueden generalizar mejor; con red fragmentada o muy
retardada, los métodos menos dependientes de un precio compartido
(\cref{th:vi-robusta}). El régimen candidato es otro: flota heterogénea,
capacidad cambiante, decisiones locales y restricciones mecánicas activas. Es
el régimen que hay que poner a prueba, no el que se supone dominado, y las
ventajas que se pueden medir son calidad y seguridad por mensaje, por ronda y
por milisegundo de cálculo, no «ser distribuido».
\end{observacion}
